% ======================================================================
% SECTION 4: DISCUSSION
% Three blocks: (1) FDA RTCT comparison, (2) AI baseline comparison,
% (3) significance for patient safety and efficacy.
% ======================================================================

\subsection{Comparison to the FDA RTCT Announcement}
\label{sec:disc-fda}

The FDA's 28 April 2026 announcement of two real-time clinical trial
proofs-of-concept and a forthcoming pilot program sets the regulatory
stage on which the four simulations in this paper operate
\cite{fda2026realtime}. The two FDA proofs-of-concept are AstraZeneca
TRAVERSE (Phase 2 mantle cell lymphoma) and Amgen STREAM-SCLC (Phase
1b limited-stage small cell lung carcinoma), with Paradigm Health as
the validated technical conduit between sites, sponsor, and the FDA
real-time stream. The four simulations here extend that proof-of-
concept along several axes:

\begin{itemize}[leftmargin=*]
\item \textbf{Simulation 1 (site)} implements a Paradigm Health-style
aggregator and the FDA real-time API stream across 4 sites (SITE-A,
SITE-B, SITE-C, SITE-D) and 116 robot instances at minute resolution
for 56 hours, with TRAVERSE, STREAM-SCLC, TRAVERSE-PED, and
NETWORK-OBS endpoint channels. The hour 00 facility diagram in
\S\ref{sec:results-sim1} shows all four channels UP at simulation
start, with a median 13-second FDA acknowledgement latency.

\item \textbf{Simulation 2 (site)} grounds a single-patient real-time
prediction trajectory in three adapted regulatory frameworks (21 CFR
312, 21 CFR 50, ICH E6(R3)) for patient PAT-2026-0042 across 1,120
days. Each stage of the journey produces a real-time signal stream
that the FDA RTCT pilot could ingest.

\item \textbf{Simulations 3 and 4 (sponsor)} scale the autonomous
sponsor side of the FDA framework from 24 hours (53 agents, 75
diagrams) to 168 hours (525 diagrams) with 168 GitHub commits across
7 branches. The sponsor server in Simulation 3 exposes 5 FastAPI
endpoints that mirror the FDA RTCT signal-sharing pattern, and the
local verification in Simulation 4 demonstrates that the same agent
logic can run on a 4 GB laptop.
\end{itemize}

The FDA's stated goal is continuous trials, where the hiatus between
phases can be eliminated or reduced. The four simulations here show
how a single 1M-token-context AI agent (Claude Code Opus 4.7 Max) can
hold the entire trial repository in context and emit hour-by-hour
real-time artifacts. The simulations are opportunistic extensions of
the FDA's vision, not critiques of it; each simulation produces an
artifact that a future FDA RTCT pilot program could in principle
ingest as input \cite{kawchak_2026_19176370,
kawchak_2026_19119939, kawchak_2026_19396256}.

A second comparison concerns advanced robotic capabilities and
predictive capabilities not present in the FDA's two proofs-of-
concept. AstraZeneca TRAVERSE and Amgen STREAM-SCLC are pharmacology
trials that report safety signals and endpoints in real time; neither
trial integrates physical robotic interventions at the site level.
The four simulations here pair the FDA's signal-sharing model with
the physical robotics layer of the four-site / 116-robot Simulation 1
network, which adds a class of safety signal (force feedback, gating
efficiency, robot authorization decision) that the pharmacology
trials do not generate. 

Each per-hour facility diagram from
Simulation 1 records, at minute resolution, the active robot
inventory by station and the per-station feedback metric. This
robotics-plus-prediction integration is a meaningful capability
extension over the pharmacology proof-of-concept, and is the area
where the four simulations here provide the most direct value-add to
the FDA's continuous-trial path.

\clearpage

\begin{table}[H]
\centering
\caption{Simulation extensions over the FDA RTCT 28 April 2026
proof-of-concept program.}
\label{tab:disc-fda-extension}
{\renewcommand{\arraystretch}{1.15}
\begin{tabular}{>{\raggedright\arraybackslash}p{3.25cm} >{\raggedright\arraybackslash}p{5.75cm} >{\raggedright\arraybackslash}p{6.3cm}}
\toprule
\textbf{Dimension} & \textbf{FDA RTCT 2026 PoC} & \textbf{Four Simulations Here} \\
\midrule
Trials                 & 2 (TRAVERSE, STREAM-SCLC)                & 4 author simulations of site, sponsor       \\
Sites                  & MD Anderson, Penn (TRAVERSE)             & 4 Simulation 1 sites, others distributed       \\
Robots                 & Not specified in the announcement         & 116 robot instances in Simulation 1                   \\
Real-time signals      & Endpoint and safety signals only         & Endpoint, safety, robot authorization, agent decision \\
Cadence                & Real-time as defined per trial            & Per-hour minute-resolution + per-stage       \\
Predictive layer       & Per-trial supervised endpoints            & Repository-scale 1M-token context across 4 simulations\\
Local-side compute     & Not addressed                            & Verified on Core i5-6200U / 4 GB     \\
\bottomrule
\end{tabular}}
\end{table}

The added predictive capability in the four simulations (1M token
context that holds protocol, EHR, robotic telemetry, and regulatory
adaptations in a single inference pass) is also crucial for patient
safety and efficacy. The FDA RTCT real-time signal-sharing channel
expects the sponsor and site to generate predictive content that the
FDA can act on; the four simulations here generate that predictive
content at a granularity (hourly per-patient narrative for Simulation
1; per-stage regulatory mapping for Simulation 2; per-hour 53-agent
decision tree for Simulations 3 and 4) that supervised AUROC models
do not produce.

\subsection{Comparison to Current AI Prediction Methods}
\label{sec:disc-ai-baseline}

The supervised baselines named in \S\ref{sec:intro-ai-baseline} differ
from the four-simulation set along several dimensions, all
computational and contextual rather than head-to-head AUROC. The
distinction is important: nothing in this paper has been benchmarked
on a fixed prediction task against a fixed test cohort under TRIPOD+AI
or CREMLS reporting standards \cite{Collins2024TRIPODAI,
ElEmam2024CREMLS}. The comparison below is scoped to the
computational signature of repository-scale processing.

\textbf{Manz 2020 versus Simulation 1.} Manz 2020 emits a single
180-day mortality probability per patient per visit, derived from a
fixed feature set, with AUC 0.89 in a 24,582-patient cohort
\cite{Manz2020Oncology180DayMortality}. Simulation 1 in
\S\ref{sec:results-sim1} emits a paragraph per minute per active
patient, plus three ASCII diagrams per hour, plus a four-Markdown
narrative bundle per hour. The Manz 2020 model cannot ingest the
robot-by-minute force-feedback log that PAT-CONT-0001's COBOT-01
biopsy generated at 1.4 N steady force, because the model is
restricted to its training feature set. The simulation here can
ingest that log and emit a narrative paragraph that mentions it by
name. The comparison is one of expressive surface area, not AUROC.

\textbf{SHIELD-RT versus Simulation 1.} SHIELD-RT is the gold-
standard prospective randomized trial for AI-directed clinical
evaluation during radiotherapy \cite{Hong2020SHIELDRT,
Elia2025SHIELDRTExternalValidation}. The original SHIELD-RT trial
covered a single radiotherapy program; Simulation 1 covers 116 robot
instances across 4 sites and 56 hours of continuous operation.
Simulation 4's 1,336-patient 168-hour cumulative throughput exceeds
the SHIELD-RT cohort scale by an order of magnitude, although the
comparison is unfair in the other direction: SHIELD-RT is real-
patient and prospectively randomized, while Simulation 4 is a
synthetic agent-rule simulation. The point of comparison is that
LLM-scale repository context can in principle scale a SHIELD-RT-like
program from one institution to many.

\textbf{SCORPIO and PROGPATH versus the four-simulation set.} SCORPIO
emits a per-patient response probability for immune checkpoint
inhibitor efficacy from routine blood tests and clinical data
\cite{Yoo2025SCORPIO}. PROGPATH emits a per-patient C-index score
across 17 external pancancer cohorts \cite{Yuan2025PROGPATH}. Both
models operate inside a single modality (blood plus clinical for
SCORPIO; pathology images for PROGPATH). The four simulations operate
across modalities (text protocol, EHR JSON, ASCII facility diagrams,
Python agent code, and regulatory citations) in the same inference
pass. The cross-modal repository context is the qualitative
extension; the four simulations cannot replace SCORPIO or PROGPATH on
their narrow tasks, but they can hold both models' outputs as inputs
to a higher-level reasoning step.

\textbf{Huang 2025 null result versus the four-simulation set.}
Huang 2025 reports that machine learning provides no significant gain
over Cox regression on real-world structured survival data
(standardized mean difference 0.01, 95\% CI -0.01 to 0.03)
\cite{Huang2025CancerSurvivalMetaAnalysis}. The four simulations do
not contradict this null result. The simulations consume an entirely
different class of input - unstructured protocol text plus telemetry
plus ASCII diagrams plus Python code - that the supervised baselines
in Huang 2025 cannot ingest at all. The null result of ML versus Cox
on structured tabular features is consistent with the observation
that the gain from LLM-scale processing comes from context that
structured-feature models do not see, not from any improvement in the
underlying regression on the structured features themselves.

\begin{table}[H]
\centering
\caption{Comparative computational signatures of supervised baselines
versus the four-simulation set. AUROC is comparable only within a
fixed task; the comparison below is scoped to the computational
expressive surface, not to head-to-head benchmark numbers.}
\label{tab:disc-baseline-vs-sim}
{\renewcommand{\arraystretch}{1.15}
\begin{tabular}{>{\raggedright\arraybackslash}p{2.6cm} >{\raggedright\arraybackslash}p{3cm} >{\raggedright\arraybackslash}p{3cm} >{\raggedright\arraybackslash}p{3cm} >{\raggedright\arraybackslash}p{3cm}}
\toprule
\textbf{Property} & \textbf{SHIELD-RT, SCORPIO} & \textbf{PROGPATH, AIM-LCpro} & \textbf{Huang 2025 Cox vs ML} & \textbf{Four Simulations} \\
\midrule
Inputs              & Fixed feature set per task    & Multimodal images plus tabular  & Tabular features only         & Repository plus narrative plus code        \\
Output              & One probability per patient   & C-index per cohort              & Hazard ratio                  & Per-hour narrative, diagrams, JSON \\
Cadence             & Per visit (months)            & Per cohort                      & Retrospective                 & Per hour, per stage                      \\
TRIPOD+AI compliant & Some                          & Some                            & Yes (meta-analytic)           & No (synthetic)                             \\
Modalities          & 1                             & 2-4                             & 1 (tabular)                   & Text plus code plus ASCII plus JSON         \\
Real-time stream    & No                            & No                              & No                            & Yes (hourly commit)                        \\
\bottomrule
\end{tabular}}
\end{table}

The four simulations in this paper do not satisfy TRIPOD+AI or CREMLS
reporting standards because they were not evaluated on real patients
\cite{Collins2024TRIPODAI, ElEmam2024CREMLS}. Any future extension of
this work toward deployment must adopt those reporting floors before
any clinical claim is made. The discussion here therefore frames the
simulations as a computational and contextual demonstration, not as a
clinical demonstration.

\subsection{Significance for Patient Safety and Efficacy Prediction}
\label{sec:disc-significance}

The central significance of the four-simulation set is the
demonstration that:

\begin{enumerate}[leftmargin=*]
\item LLM-scale repository context (1M tokens) produces patient
prediction artifacts at a cadence (1 hour) and depth (3 ASCII
diagrams plus 4 Markdown files per hour for Simulation 1; 24 Python
scripts plus 75 diagrams plus JSON per 24 hours for Simulation 3;
525 diagrams plus 168 scripts plus 168 JSON per 168 hours for
Simulation 4) that no current narrow supervised oncology model
matches.



\item The same repository-scale context can be operated by a single
autonomous agent: a sponsor in Simulations 3 and 4, a site
coordinator in Simulation 1, a single-patient orchestrator in
Simulation 2. This opens a path to fully autonomous safety and
efficacy prediction loops aligned with the FDA RTCT pilot program.

\item Simulation 4's local verification on a Core i5-6200U with 4 GB
RAM proves that the computational floor is far below typical data-
center infrastructure - a result that strengthens the FDA's
observation that real-time trials are \emph{"not only possible, but
also potentially transformative"} \cite{fda2026realtime}.
\end{enumerate}

The work is opportunistic with respect to the FDA's framework: it
offers concrete file-level artifacts that the RTCT pilot program could
ingest as inputs once real-patient validation is complete. None of
the four simulations claims clinical deployment readiness; the value
is computational and contextual. Cite \cite{fda2026realtime,
kawchak_2026_19176370, kawchak_2026_19244918, kawchak_2026_18810541,
kawchak_2026_19119939, kawchak_2026_19396256} for the underlying
artifacts.

\subsection{Cloud Versus Local Compute Trade-Offs}
\label{sec:disc-cloud-local}

A practical question for any RTCT-aligned deployment is whether the
sponsor-side AI runs on cloud compute or on local hardware. The four
simulations here cover both modes. Simulations 1, 2, and 3 ran on
cloud compute; Simulation 4 ran on cloud compute and was verified
locally on a 4 GB Core i5-6200U laptop. The trade-offs are
summarized in the diagram below.

\begin{verbatim}
+============================================================================+
|          CLOUD VS LOCAL COMPUTE FOR RTCT-ALIGNED ONCOLOGY TRIAL AI         |
+----------------------+--------------------+--------------------------------+
| Property             | Cloud Compute Only | Cloud + Local Verification     |
+----------------------+--------------------+--------------------------------+
| Wall-clock speed     | Highest            | Cloud high, local slower       |
| 1M token context     | Yes (Claude Opus)  | Yes (cloud) / partial (local)  |
| Reproducibility      | Cloud commit log   | Cloud + local rerun parity     |
| Hardware floor       | Data-center scale  | i5-6200U / 4 GB demonstrated   |
| PHI / data security  | Cloud transit risk | Local processing + cloud meta  |
| Audit trail          | GitHub commits     | GitHub commits + local hash    |
| Site adoption cost   | Cloud subscription | Off-the-shelf laptop           |
+----------------------+--------------------+--------------------------------+
\end{verbatim}

The advantage of cloud compute is that a single 1M-token-context
inference pass can hold the entire trial repository in memory and
emit per-hour artifacts in seconds; the disadvantage is that PHI
traveling to and from the cloud must be carefully governed, and that
the audit trail depends on the cloud provider's transparency. The
advantage of cloud-plus-local verification is that the local rerun
can validate the cloud output on hardware that an FDA inspector or a
trial sponsor's IT team can audit physically; the disadvantage is
that local hardware introduces operating-system-specific quirks and
cannot itself host a 1M-token-context inference pass.

The optimal future configuration combines both. A future Claude Code
or competing AI local instance that runs simultaneously on cloud
(for the 1M-token context full-repository reasoning) and on local
hardware (for per-task specialist agents that handle PHI without
leaving the site) would offer both the power of cloud compute and
the security and flexibility of local computing. Such a configuration
is the natural endpoint of the trade-off table above and is the
subject of Track B in the Future Work section
(\S\ref{sec:limitations}).

\subsection{Code-Based Versus Text-Only Simulations - Practical Implications}
\label{sec:disc-code-vs-text}

A second practical question is whether the sponsor-side AI emits
Python code or only Markdown plus ASCII text. Simulation 1 emits
text only; Simulations 2, 3, and 4 emit Python plus JSON. The
practical trade-offs were sketched in Methods
(\S\ref{sec:methods-code-vs-text}); the implications for performance
in real RTCT-aligned deployments deserve a fuller discussion here.

Code-based simulations win on \emph{practicality} for downstream
automation. A Python script can re-execute on local hardware to
verify the cloud-generated output (Simulation 4 demonstrates this on
a 4 GB laptop), an importer can call the agent modules directly, and
a CI workflow can lint and test the code under the repository's
\texttt{ruff} configuration. Text-only simulations win on
\emph{auditability} for clinical readers. A radiation oncologist
reviewing Simulation 1's hour 12 patient flow diagram can read the
PAT-CONT-0051A through PAT-CONT-0058A bracketed timeline directly
without running any code, which lowers the technical barrier to
clinical review.

The performance trade-off in cloud compute is also asymmetric.
Generating Markdown and ASCII is computationally cheap on the cloud
side because the model only needs to produce text tokens; generating
Python code is more expensive because the model must produce
executable structure under the constraint that the code parses and
runs. Simulation 1 generated 56 hours of text in roughly 4 to 5
hours of cloud time, whereas Simulation 4 generated 168 hours of
Python plus JSON in roughly 12 to 14 hours of cloud time. The
disparity reflects the higher per-token cost of code generation, not
any difference in the underlying 1M-token context capability.

For a real RTCT-aligned deployment, the practical recommendation is
to combine text-only artifacts (for clinical readers) and code
artifacts (for downstream automation) in a single trial. Simulation 1
plus Simulation 4 paired together is the closest existing
approximation: Simulation 1 generates the per-hour patient narrative
that a clinical reader can audit, and Simulation 4 generates the per-
hour sponsor decision script that a CI workflow can re-run. The
synergistic value of the pair exceeds the value of either alone.
